\section{Nine positive unknowns: a second improved equation}
\label{sec:nine}
The same support-selective source also improves the nine-variable Jones
route.  Its primary circuit has nine strictly positive unknowns, exact
total degree $19\,615\,180$, and $606$ binary arithmetic operations.  The
equivalent version over nonnegative integers uses $609$ operations.  The
previous package-45 positive equation has degree $129\,875\,980$ and $788$
operations.  This is a reduction of approximately $84.9\%$ in degree and
$23.1\%$ in operations, without changing the number of unknowns.

Let the raw quadratic source, its weights and its targets be as above, and
write $W=175$, $H=5148$, and $K=H+2W=5498$.  Set
$M=2^{2140017}$ and $b_0=2^{2140030}=8192M$.  Initially take
$\epsilon,\phi,n\geq0$ and $g,h,s,w,i,j>0$. Define
\begin{align*}
b&=b_0+\epsilon,& B&=b^3,& \ell&=\sum_{r=1}^{31}B^{w_r},\\
N_0&=B^{W+1},&N_1&=4B^{K+1},&N&=N_0N_1,\\
L&=N_0-1-(b-1)\ell,& Z&=\sum_{r=1}^{22}B^{E_r},\\
T&=L+(B-2)ZN_0,&&&\\
J_0&=2(1+\ell+g)^2\Dc(B)+B\sum_{r=0}^{K}B^r,&&&\\
S&=g+J_0N_0,&R&=(N-1)((S+T+1)N+T+1),\\
G_{\mathrm{cap}}&=32bB^W.&&&
\end{align*}
Every exponent is fixed.  The raw residuals are evaluated at positive
digits $x_r=z_r+1$: the one added $\ell$ implements all thirty-one
positivity shifts in the common convolution factor.

The coefficient at the largest exponent of $\Dc$ is $2$, and all its other
coefficients have absolute value below $M$. Hence $\Dc(B)>0$ unconditionally.
The source therefore satisfies $R,N\geq8$ and $b\leq R,N$ before any
equation is imposed. If $g<4bB^W$, then $1+\ell+g<5bB^W$, and the choice
$b\geq8192M$ gives $0<S,T<N$.  With $b$ a power of two, the exact
population-count identity
\[
\operatorname{popcount}(R)=2\log_2N+
\operatorname{popcount}(S+T)-\operatorname{popcount}(S)
-\operatorname{popcount}(T)
\]
shows that $N^2\mid\binom{2R}{R}$ is precisely the no-carry condition.
Its low block encodes $g=\sum z_rB^{w_r}$ with $0\leq z_r<b$; its target
blocks impose all the original positive quadratic equations simultaneously.
Conversely, every positive parent solution permits arbitrarily large
dyadic bases with $32g<G_{\mathrm{cap}}$.  The all-one parent assignment
is not a zero, because its residual $a-E-sN$ equals $-1$, so $g>0$.

All the numerical Jones bridge registers are as follows
\cite{package45,jones}. The names $K_1,D_1,H_1$ distinguish these
registers from the source exponent bound $K$, kernel $\Dc$, and maximum
target $H$:
\begin{align*}
U&=N^2w,&Y_0&=N^2s,&M_0&=RY_0,\\
P&=2M_0^2U,&K_1&=R+1+h(P-1),&A&=M_0(U+1),\\
B_1&=2R+1,&W_1&=bw,&C&=B_1+W_1+\phi,\\
J&=4A-5,&Q&=3W_1C-2(W_1^2-1),&&\\
D_1&=(A^2-1)C^2+1,&E&=iJD_1C^2,&&\\
F&=(A^2-1)E^2+1,&G&=A+F(CD_1+1-A),&&\\
H_1&=B_1+jC,&I&=(G^2-1)H_1^2+1,&&\\
a_0&=(P^2-1)K_1^2+1,&a_1&=D_1FI,&&\\
u&=JF,&v&=QF+J(H_1-C).&&
\end{align*}
Define the margin polynomial
\[
m=K_1^2G_{\mathrm{cap}}
 -4(C-K_1Y_0)^2G_{\mathrm{cap}}-8gK_1^2.
\]
Here $m$ is a polynomial register, not a tenth unknown.

The hypotheses of \texttt{Jones1982.Nine.ratio\_sound} and
\texttt{ratio\_complete} are $R,N\geq8$, $b\leq R,N$,
$b,G_{\mathrm{cap}}>0$, $h,s,w,i,j>0$ and $\phi\geq0$.
All have been supplied independently of the carry condition.
Soundness states that $a_0,a_1$ squares, $u\mid v$, and $m>0$ imply
that $b$ is a power of two, $N^2\mid\binom{2R}{R}$, and
$8g<G_{\mathrm{cap}}$. The latter is exactly $g<4bB^W$, the
interval needed by the source. Completeness starts with $b$ a power
of two, the divisibility, and $32g<G_{\mathrm{cap}}$, and supplies
$h,s,w,i,j,\phi$ satisfying the four arithmetic conditions.
This is precisely the stronger interval supplied by a parent zero.
No property of the Sun ten-integer bridge is used here.

Set
\[
V_0=PK_1,\qquad V_1=(AC)(AE)(GH_1),\qquad M_d=Q+H_1.
\]
These registers are positive. The identities
\begin{align*}
V_0^2-a_0&=K_1^2-1>0,&(AC)^2-D_1&=C^2-1>0,\\
(AE)^2-F&=E^2-1>0,&(GH_1)^2-I&=H_1^2-1>0
\end{align*}
give $0<a_0<V_0^2$ and $0<a_1<V_1^2$ on every outer assignment.
Also $Q>0$, $H_1-C>0$, and
$0<v/u=Q/J+(H_1-C)/F<M_d-1$.

One polynomial combines these four conditions. Put
\[
\Omega=um,\qquad
\Theta=u(n+M_d)+v-\Omega(2M_d+V_0V_1),
\]
and take
\[
\boxed{\bigl(\Theta^2-\Omega^2(a_0+V_0^2a_1)\bigr)^2
       -4\Omega^4V_0^2a_0a_1=0.}
\]
At $m=0$ its left side is nonzero. Otherwise its four-factor norm identity
implies that a rational number is $r_0+V_0r_1$, with $r_t^2=a_t$.
The strict majorant excludes a zero rational sum; the identity
$2(r_0+V_0r_1)r_0=(r_0+V_0r_1)^2+a_0-V_0^2a_1$ then makes both roots
integers. Their bounds give
$2M_d+V_0V_1+r_0+V_0r_1\geq2M_d+1>0$, which forces $m>0$ and
$u\mid v$.  The converse chooses $n$ from this same positive expression.
Thus this finalizer has both the required soundness and completeness.

To define the polynomial on nine \emph{positive} coordinates, replace
$\epsilon,\phi,n$ throughout by their positive coordinates minus one,
leaving $g,h,s,w,i,j$ positive. For the nine \emph{nonnegative} version,
instead leave $\epsilon,\phi,n$ unchanged and replace each of
$g,h,s,w,i,j$ by its nonnegative coordinate plus one. These are bijections
with the mixed domain used in the proof. They establish the stated
solvability equivalences; they do not justify an unrestricted-integer
interpretation of this nine-variable equation.

The full exact degree is
\[
1152t+12W+280=19\,615\,180,
\qquad t=3(H+3W+2)=17\,025.
\]
The cancellation-resolved circuit proves this upper bound, and its
all-variable diagonal has nonzero leading coefficient $2^{2050}$, proving
equality. Its critical exact identity is
\begin{align*}
\Delta_V=V_1^2-a_1
 ={}&(C^2-1)(AE)^2(GH_1)^2\\
   &+D_1(E^2-1)(GH_1)^2+D_1F(H_1^2-1).
\end{align*}
Writing $E_0=u(n+M_d)+v-2\Omega M_d$, the inner polynomial of the
final equation equals
\[
E_0^2-2E_0\Omega V_0V_1+\Omega^2V_0^2\Delta_V-\Omega^2a_0.
\]
Its uniquely highest degree is $576t+6W+140$, with diagonal coefficient
$2^{1025}$. The final subtracted term has smaller degree than its square.
The full proof is in \file{nine/PROOF.md}; the polynomial identity is
independently replayed by the verifier. The source is explicit, while the
finite audits are implementation checks, not a decision of consistency
or a new end-to-end formal proof of the inherited compiler.
